\section{Part II.a (ii): Deflection, Precession, and Gravitational Waves}

\subsection{Deflection}

The isotropic sector — the degenerate locus of A4 with a single scalar scale — produces exactly half the observed light-bending effect around a massive body. The effect depends on a single scalar, the gravitational potential $\Phi$, and at grazing incidence the prediction is a factor of two below the measured value. This file establishes that the missing half comes from spatial scale anisotropy, which a scalar scale cannot carry. The factor of two is thus not a fitting uncertainty but a structural statement: a theory with one scale can only supply the temporal contribution. Numerical verification is given for a ray grazing the Sun, matching VLBI measurements.

\begin{definition}[\lean{Deflection.deflectionIso}]
The grazing deflection produced by the temporal scale alone — the whole of what an isotropic (single-scalar) scale field can produce:
\[
\text{deflectionIso}(G, M, b, c) = \frac{2GM}{bc^2}.
\]
\end{definition}

\begin{definition}[\lean{Deflection.deflectionObs}]
The measured grazing deflection:
\[
\text{deflectionObs}(G, M, b, c) = \frac{4GM}{bc^2}.
\]
\end{definition}

\begin{theorem}[\lean{Deflection.iso\_is\_exactly\_half}]
The isotropic sector gives exactly half the observed deflection, with no adjustable quantity anywhere in the statement:
\[
2 \cdot \text{deflectionIso}(G, M, b, c) = \text{deflectionObs}(G, M, b, c).
\]
\begin{proof}
Immediate from the definitions by ring algebra.
\end{proof}
\end{theorem}

\begin{theorem}[\lean{Deflection.iso\_ratio\_half}]
Stated as a ratio, where the denominator is nonzero. If $\text{deflectionObs}(G, M, b, c) \neq 0$, then
\[
\frac{\text{deflectionIso}(G, M, b, c)}{\text{deflectionObs}(G, M, b, c)} = \frac{1}{2}.
\]
\begin{proof}
The hypothesis that the observed deflection is nonzero implies the isotropic deflection is nonzero. By the previous theorem, the observed deflection is twice the isotropic one. Dividing yields the result.
\end{proof}
\end{theorem}

\begin{theorem}[\lean{Deflection.missing\_half}]
The deficit is exactly the isotropic prediction over again: the missing contribution is the same size as the one a scalar scale can produce. In the scale reading, the temporal and the spatial scale contribute equally, and a single scalar carries only the temporal one:
\[
\text{deflectionObs}(G, M, b, c) - \text{deflectionIso}(G, M, b, c) = \text{deflectionIso}(G, M, b, c).
\]
\begin{proof}
Immediate from the definitions by ring algebra.
\end{proof}
\end{theorem}

\begin{definition}[\lean{Deflection.solarObsArcsec}]
The measured-branch (``observed'') prediction for grazing deflection of a ray around the Sun, in arcseconds:
\[
\text{solarObsArcsec} = \frac{4 \times 6.674 \times 10^{-11} \times 1.989 \times 10^{30}}{6.957 \times 10^8 \times (2.998 \times 10^8)^2} \times 206264.806.
\]
\end{definition}

\begin{definition}[\lean{Deflection.solarIsoArcsec}]
The isotropic-sector prediction for the same grazing deflection, in arcseconds:
\[
\text{solarIsoArcsec} = \frac{2 \times 6.674 \times 10^{-11} \times 1.989 \times 10^{30}}{6.957 \times 10^8 \times (2.998 \times 10^8)^2} \times 206264.806.
\]
\end{definition}

\begin{theorem}[\lean{Deflection.solarObs\_value}]
The observed prediction yields $1.7515 \text{\textquotedbl}$, matching the VLBI value $1.7509 \pm 0.0002 \text{\textquotedbl}$ to about $0.04\%$:
\[
1.7515 < \text{solarObsArcsec} < 1.7516.
\]
\begin{proof}
By numerical verification.
\end{proof}
\end{theorem}

\begin{theorem}[\lean{Deflection.solarIso\_value}]
The isotropic sector predicts $0.8758 \text{\textquotedbl}$, wrong by a factor of two. This is the sharpest quantitative statement in the development, and it is a refutation of the scalar axiom rather than a confirmation of anything:
\[
0.87 < \text{solarIsoArcsec} < 0.88.
\]
\begin{proof}
By numerical verification.
\end{proof}
\end{theorem}

\begin{theorem}[\lean{Deflection.solar\_ratio}]
The two solar numbers stand in the exact ratio $2:1$, as the general theorem requires:
\[
\text{solarObsArcsec} = 2 \cdot \text{solarIsoArcsec}.
\]
\begin{proof}
Immediate by ring algebra from the definitions.
\end{proof}
\end{theorem}

\begin{theorem}[\lean{Deflection.iso\_fails\_by\_far}]
The isotropic sector misses the measured value by more than $0.8 \text{\textquotedbl}$, which is four thousand times the VLBI uncertainty: the failure is not a modelling imprecision but a structural one:
\[
0.87 < \text{solarObsArcsec} - \text{solarIsoArcsec}.
\]
\begin{proof}
By the ratio theorem and the numerical bounds on the isotropic value, the difference exceeds $0.87 \text{\textquotedbl}$, whereas the VLBI uncertainty is $\pm 0.0002 \text{\textquotedbl}$.
\end{proof}
\end{theorem}

\subsection{PPN}

The measured deflection angle can be parametrised in terms of two dimensionless numbers: $\gamma$, weighting the spatial scale contribution, and $\beta$, measuring the nonlinearity of the temporal scale. The standard value is $\gamma = 1$. The isotropic sector, having no spatial scale at all, effectively sets $\gamma = 0$. The framework's scale reciprocity condition — the product of temporal and radial scales equals one — forces $\gamma = 1$ exactly to first order. This file proves that reciprocity determines the spatial-scale coefficient, paying the debt of the isotropic sector through a necessary relation among components.

\begin{definition}[\lean{PPN.deflectionGamma}]
Grazing deflection with the spatial scale weighted by $\gamma$:
\[
\text{deflectionGamma}(G, M, b, c, \gamma) = \left(\frac{2GM}{bc^2}\right)(1 + \gamma).
\]
Here $\gamma = 0$ keeps only the temporal scale, and $\gamma = 1$ is the measured case.
\end{definition}

\begin{theorem}[\lean{PPN.gamma\_zero\_is\_isotropic}]
With no spatial contribution the formula returns the isotropic sector's prediction — the half that a scalar scale can produce:
\[
\text{deflectionGamma}(G, M, b, c, 0) = \text{deflectionIso}(G, M, b, c).
\]
\begin{proof}
Immediate from the definitions by ring algebra.
\end{proof}
\end{theorem}

\begin{theorem}[\lean{PPN.gamma\_one\_is\_observed}]
With the spatial scale weighted equally, the formula returns the measured value:
\[
\text{deflectionGamma}(G, M, b, c, 1) = \text{deflectionObs}(G, M, b, c).
\]
\begin{proof}
Immediate from the definitions by ring algebra.
\end{proof}
\end{theorem}

\begin{theorem}[\lean{PPN.reciprocal\_expansion}]
If the radial factor is the exact inverse of the temporal one — which is what the reciprocity relation $s_t \cdot s_r = 1$ says — then the expansion identity holds:
\[
(1 - 2\mu)^{-1} = 1 + 2\mu + \frac{4\mu^2}{1 - 2\mu},
\]
where $1 - 2\mu \neq 0$. This is an identity, not an approximation. The coefficient of $\mu$ is $2$, and the remainder is displayed rather than dropped.
\begin{proof}
Clear the denominators and verify by field algebra.
\end{proof}
\end{theorem}

\begin{theorem}[\lean{PPN.gamma\_eq\_one\_of\_reciprocal}]
Reciprocity forces $\gamma = 1$. Matching $g_{rr} = 1 + 2\gamma\mu$ against the inverse of $g_{tt} = 1 - 2\mu$ at first order in $\mu$, the reciprocal condition leaves no freedom. Stated exactly: if $\mu \neq 0$, $1 - 2\mu \neq 0$, and the radial factor equals $1 + 2\gamma\mu$ and also equals $(1-2\mu)^{-1}$, then
\[
\gamma = 1 + \frac{2\mu}{1 - 2\mu}.
\]
\begin{proof}
By the reciprocal expansion theorem, $(1-2\mu)^{-1} = 1 + 2\mu + \frac{4\mu^2}{1-2\mu}$. Equating this to $1 + 2\gamma\mu$ gives $2\gamma\mu = 2\mu + \frac{4\mu^2}{1-2\mu}$. Dividing by the nonzero quantity $2\mu$ and clearing denominators yields the result.
\end{proof}
\end{theorem}

\begin{theorem}[\lean{PPN.gamma\_eq\_one\_first\_order}]
In the strict first-order regime — where $\mu^2$ is discarded, exactly as in the Newtonian limit — reciprocity gives $\gamma = 1$ on the nose. If $\mu \neq 0$ and $1 + 2\gamma\mu = 1 + 2\mu$, then $\gamma = 1$.
\begin{proof}
The hypothesis gives $2\gamma\mu = 2\mu$. Since $\mu \neq 0$, divide by $2\mu$ to obtain $\gamma = 1$.
\end{proof}
\end{theorem}

\begin{theorem}[\lean{PPN.deflection\_reciprocal\_eq\_observed}]
The debt is paid at the level of the coefficient. If $\gamma = 1$, then
\[
\text{deflectionGamma}(G, M, b, c, \gamma) = \text{deflectionObs}(G, M, b, c).
\]
The half that the isotropic sector could not produce is the radial scale, and its size was never free.
\begin{proof}
Substitute $\gamma = 1$ and apply the theorem that $\gamma = 1$ returns the observed value.
\end{proof}
\end{theorem}

\begin{theorem}[\lean{PPN.halves\_equal}]
The two halves are equal in size. The contribution from the spatial scale ($\gamma = 1$ term) minus the contribution from the temporal scale alone ($\gamma = 0$ term) equals the temporal contribution:
\[
\text{deflectionGamma}(G, M, b, c, 1) - \text{deflectionGamma}(G, M, b, c, 0) = \text{deflectionGamma}(G, M, b, c, 0).
\]
\begin{proof}
Immediate from the definitions by ring algebra.
\end{proof}
\end{theorem}

\begin{theorem}[\lean{PPN.solar\_reciprocal\_value}]
The reciprocal ($\gamma = 1$) prediction for a ray grazing the Sun, in arcseconds, is $1.7515 \text{\textquotedbl}$, against the VLBI value $1.7509 \pm 0.0002 \text{\textquotedbl}$:
\[
1.7515 < \text{solarObsArcsec} < 1.7516.
\]
\begin{proof}
This is the same bound as the observed deflection, since reciprocity forces $\gamma = 1$ and thus agreement with the measured value.
\end{proof}
\end{theorem}

\begin{theorem}[\lean{PPN.gamma\_prediction\_is\_sharp}]
Current solar-system bounds from the Cassini tracking data put $\gamma$ within about $2 \times 10^{-5}$ of $1$. The framework's reciprocity condition predicts $\gamma = 1$ with no free parameter. If $\gamma = 1$, then
\[
\gamma - 1 = 0.
\]
\begin{proof}
Immediate from the hypothesis.
\end{proof}
\end{theorem}

\subsection{Precession}

Mercury's perihelion advance per orbit is the second solar-system test, of a different kind and with a different failure factor in the isotropic sector. The coefficient in the standard parametrisation is $(2 - \beta + 2\gamma)/3$, where $\gamma$ is the spatial-scale weight (forced to $1$ by reciprocity) and $\beta$ measures the nonlinearity of the temporal scale (equals $1$ from A6′). With both parameters at their framework values, the coefficient is $1$ and the prediction is measured. The isotropic sector, having no spatial scale, gives $1/3$. The two tests fail by factors $2$ and $3$ respectively, which are different and thus cannot both be absorbed by rescaling. Both arise from the same single parameter, indicating a single missing ingredient: the spatial scale.

\begin{definition}[\lean{Precession.coeff}]
The perihelion coefficient in terms of the temporal nonlinearity $\beta$ and the spatial-scale weight $\gamma$:
\[
\text{coeff}(\beta, \gamma) = \frac{2 - \beta + 2\gamma}{3}.
\]
\end{definition}

\begin{theorem}[\lean{Precession.coeff\_reciprocal}]
The reciprocal condition gives the measured coefficient. Here $\gamma = 1$ is forced by $s_t \cdot s_r = 1$ and $\beta = 1$ is the scale response of A6′:
\[
\text{coeff}(1, 1) = 1.
\]
\begin{proof}
By direct computation.
\end{proof}
\end{theorem}

\begin{theorem}[\lean{Precession.coeff\_isotropic}]
The isotropic sector gives one third, since with no spatial scale $\gamma = 0$:
\[
\text{coeff}(1, 0) = \frac{1}{3}.
\]
\begin{proof}
By direct computation.
\end{proof}
\end{theorem}

\begin{theorem}[\lean{Precession.isotropic\_is\_third}]
The two coefficients differ by exactly a factor of three — a different factor from the two of light deflection, obtained from the same single parameter:
\[
\text{coeff}(1, 0) \times 3 = \text{coeff}(1, 1).
\]
\begin{proof}
By direct computation from the coefficients.
\end{proof}
\end{theorem}

\begin{theorem}[\lean{Precession.coeff\_affine\_in\_gamma}]
The sensitivity to $\gamma$ is what does the work: the coefficient is affine in $\gamma$ with slope $2/3$:
\[
\text{coeff}(\beta, \gamma) = \text{coeff}(\beta, 0) + \frac{2}{3}\gamma.
\]
\begin{proof}
Expand the definition and verify by ring algebra.
\end{proof}
\end{theorem}

\begin{definition}[\lean{Precession.mercuryFactor}]
Everything in the Mercury prediction except the factor of $\pi$:
\[
\text{mercuryFactor} = \frac{6 \times 6.674 \times 10^{-11} \times 1.989 \times 10^{30} \times 415.2030829 \times 206264.806}{5.7909 \times 10^{10} \times (1 - 0.2056^2) \times (2.998 \times 10^8)^2}.
\]
\end{definition}

\begin{theorem}[\lean{Precession.mercuryFactor\_bounds}]
The bounds on the Mercury factor are
\[
13.6837 < \text{mercuryFactor} < 13.6838.
\]
\begin{proof}
By numerical verification.
\end{proof}
\end{theorem}

\begin{definition}[\lean{Precession.mercuryPrecession}]
Predicted perihelion advance of Mercury, arcseconds per century, at the reciprocal value $\gamma = 1$:
\[
\text{mercuryPrecession} = \pi \times \text{mercuryFactor}.
\]
\end{definition}

\begin{theorem}[\lean{Precession.mercury\_value}]
The framework predicts $42.989 \text{\textquotedbl}$ per century, against the measured $42.98 \pm 0.04 \text{\textquotedbl}$. The bound is certified to $\pm 0.001 \text{\textquotedbl}$, which is what the comparison to a $\pm 0.04 \text{\textquotedbl}$ measurement requires:
\[
42.988 < \text{mercuryPrecession} < 42.990.
\]
\begin{proof}
From the bounds on the Mercury factor and the bounds on $\pi$, use the positivity of $\pi$ and algebraic manipulation to establish the tight bounds.
\end{proof}
\end{theorem}

\begin{definition}[\lean{Precession.mercuryIsotropic}]
The isotropic sector's prediction, one third of the full prediction:
\[
\text{mercuryIsotropic} = \frac{\text{mercuryPrecision}}{3}.
\]
\end{definition}

\begin{theorem}[\lean{Precession.mercury\_isotropic\_value}]
The isotropic sector predicts $14.33 \text{\textquotedbl}$ — wrong by nearly $29 \text{\textquotedbl}$, some seven hundred times the measurement error. As with light bending, the failure is structural:
\[
14.3 < \text{mercuryIsotropic} < 14.4.
\]
\begin{proof}
Divide the bounds on the full prediction by $3$.
\end{proof}
\end{theorem}

\begin{theorem}[\lean{Precession.two\_tests\_different\_factors}]
The two solar-system tests fail in the isotropic sector by different factors — two and three — both produced by setting $\gamma = 0$ in the same parametrisation. A single missing spatial scale accounts for both:
\[
\text{coeff}(1, 0) \times 3 = \text{coeff}(1, 1) \quad \text{and} \quad 2 \times \frac{1}{2} = 1.
\]
\begin{proof}
The first equation is the isotropic-is-third theorem. The second is immediate.
\end{proof}
\end{theorem}

\subsection{Covariance}

The framework's axiom A6 states that the universe dissipates: the fiducial scale drifts without end. The obvious objection is that this should wreck local physics — if the unit of length is running away, why is the hydrogen atom the same size today as yesterday, and why does the solar system not expand? The answer is forced by axiom A5 and is proved here. Cosmic dissipation is a spatially constant drift of the fiducial, so every geometric object built from derivatives of the scale is unaffected. Local geometry is thus exactly time-independent, and the field equation is form-invariant. Bound systems do not expand, and this is a theorem rather than an assumption. The expansion is visible only in the comparison of different fiducials — which is exactly what a redshift measurement is.

\begin{definition}[\lean{Covariance.CosmicHistory}]
A cosmic history is a data structure encoding a local scale profile together with a global fiducial drift. Specifically, it consists of:
\begin{itemize}
\item A local scale profile $\text{base}$, representing the gravitational field,
\item A global fiducial drift $\text{drift} : \mathbb{R} \to A$, produced by dissipation,
\item A proof that the drift is spatially constant: for every $t \in \mathbb{R}$ and every spatial index $i \in \text{Fin } n$, the spatial derivative $d_i(\text{drift}(t)) = 0$.
\end{itemize}
This encodes that a fiducial is by definition the same everywhere at a given epoch, so only the uniform scaling changes in time.
\end{definition}

\begin{theorem}[\lean{Covariance.CosmicHistory.sig\_eq}]
The scale gradient does not drift. For a cosmic history $H$ and any cosmic time $t$ and spatial index $i$,
\[
\text{sig}(H.\text{scale}(t))(i) = \text{sig}(H.\text{base})(i).
\]
\begin{proof}
The scale field at time $t$ is the base plus a spatially-constant drift. When taking spatial derivatives, the constant part vanishes.
\end{proof}
\end{theorem}

\begin{theorem}[\lean{Covariance.CosmicHistory.Chr\_eq}]
The connection is epoch-independent. For any cosmic times $t$ and spatial indices $a, b, e$,
\[
\text{Chr}(H.\text{scale}(t))(a, b, e) = \text{Chr}(H.\text{base})(a, b, e).
\]
\begin{proof}
The Christoffel symbols depend only on derivatives of the scale. Since the drift is spatially constant, these derivatives are the same at all epochs.
\end{proof}
\end{theorem}

\begin{theorem}[\lean{Covariance.CosmicHistory.Rm\_eq}]
Curvature is epoch-independent. Two observers at cosmic times $t$ and $t'$, however far apart, measure exactly the same local geometry. For any indices $a, b, e, f \in \text{Fin } n$,
\[
\text{Rm}(H.\text{scale}(t))(a, b, e, f) = \text{Rm}(H.\text{scale}(t'))(a, b, e, f).
\]
\begin{proof}
Both the connection and curvature are unchanged when the scale is shifted by a spatially constant amount.
\end{proof}
\end{theorem}

\begin{theorem}[\lean{Covariance.CosmicHistory.Ric\_eq\_of\_drift}]
The Ricci tensor is epoch-independent. For any cosmic times $t, t'$ and spatial indices $b, e$,
\[
\text{Ric}(H.\text{scale}(t))(b, e) = \text{Ric}(H.\text{scale}(t'))(b, e).
\]
\begin{proof}
The Ricci tensor is a contraction of the Riemann tensor, which is epoch-independent.
\end{proof}
\end{theorem}

\begin{theorem}[\lean{Covariance.CosmicHistory.Defm\_eq}]
The scale deformation — the source term of the field equation — is epoch-independent. For any cosmic times $t, t'$ and spatial indices $i, j$,
\[
\text{Defm}(n, H.\text{scale}(t))(i, j) = \text{Defm}(n, H.\text{scale}(t'))(i, j).
\]
\begin{proof}
The deformation tensor depends on second derivatives of the scale, which are unchanged by a spatially constant shift.
\end{proof}
\end{theorem}

\begin{theorem}[\lean{Covariance.CosmicHistory.field\_equation\_form\_invariant}]
The field equation is form-invariant under cosmic dissipation. If $e^{2\sigma}\mathcal{R} = \kappa\rho$ holds at one epoch, it holds at every epoch with the same $\kappa$ and the same local $\rho$. Global dissipation therefore does not renormalise local gravity, and local general covariance is preserved exactly. For any $\kappa, \rho \in A$ and cosmic times $t, t'$, if
\[
\text{RscBare}(n, H.\text{scale}(t)) = \kappa \rho,
\]
then
\[
\text{RscBare}(n, H.\text{scale}(t')) = \kappa \rho.
\]
\begin{proof}
The scalar curvature is a sum of Ricci components, each of which is epoch-independent.
\end{proof}
\end{theorem}

\begin{theorem}[\lean{Covariance.CosmicHistory.no\_local\_expansion}]
Bound systems do not expand. Collected statement: everything locally measurable is exactly constant along the cosmic drift. Expansion is not a local force that space exerts on matter; it is only visible when two different fiducials are compared, which is what a cosmological redshift measurement does. For any cosmic times $t, t'$:
\[
\begin{aligned}
&\left(\forall i \in \text{Fin } n, \, \text{sig}(H.\text{scale}(t))(i) = \text{sig}(H.\text{scale}(t'))(i)\right) \\
&\wedge \left(\forall a, b, e, f \in \text{Fin } n, \, \text{Rm}(H.\text{scale}(t))(a,b,e,f) = \text{Rm}(H.\text{scale}(t'))(a,b,e,f)\right) \\
&\wedge \left(\forall b, e \in \text{Fin } n, \, \text{Ric}(H.\text{scale}(t))(b, e) = \text{Ric}(H.\text{scale}(t'))(b, e)\right) \\
&\wedge \left(\forall i, j \in \text{Fin } n, \, \text{Defm}(n, H.\text{scale}(t))(i, j) = \text{Defm}(n, H.\text{scale}(t'))(i, j)\right) \\
&\wedge \, \text{RscBare}(n, H.\text{scale}(t)) = \text{RscBare}(n, H.\text{scale}(t')).
\end{aligned}
\]
\begin{proof}
Each component follows from the epoch-independence theorems proved above.
\end{proof}
\end{theorem}

\subsection{Singularity}

In general relativity a singularity is where the metric degenerates and curvature invariants blow up. In this framework that cannot happen, and the reason is structural rather than a matter of finding better solutions. Axiom A2 makes the scale $s = e^\sigma$ an element of the unit group $A^\times$. A unit is invertible by definition. This removes singularities entirely: the conformal factor never vanishes, the energy scale is likewise always defined, and curvature is a polynomial in derivatives of $\sigma$ in which $s$ does not appear — so no divergence can be produced by the scale being small. What general relativity reports as a singularity is, here, a region where $\sigma$ becomes very negative: the scale is tiny and the energy scale is enormous, but both are finite and invertible. The theory keeps speaking.

\begin{theorem}[\lean{Singularity.ScaleField.scale\_ne\_zero}]
The scale never vanishes. A unit of measure is by construction something one can divide by. For a scale field $F$ with $A$ nontrivial,
\[
(F.s : A) \neq 0.
\]
\begin{proof}
Suppose the scale vanishes. Then $F.s \cdot F.s^{-1} = 0 \cdot F.s^{-1} = 0$, but by definition of a unit in the group of units, this product equals $1$. This gives $1 = 0$, contradicting nontriviality.
\end{proof}
\end{theorem}

\begin{theorem}[\lean{Singularity.ScaleField.energy\_ne\_zero}]
The energy scale never vanishes either, by the duality A3. For a scale field $F$ with $A$ nontrivial,
\[
F.\text{en} \neq 0.
\]
\begin{proof}
By the duality relation $F.\text{en} \cdot F.s = 1$. If the energy scale vanished, the product would be zero, but it must equal $1$ by definition.
\end{proof}
\end{theorem}

\begin{theorem}[\lean{Singularity.ScaleField.conformal\_factor\_isUnit}]
The metric never degenerates. The conformal factor $e^{2\sigma}$ relating the measured line element to the bare one is a unit at every point:
\[
\text{IsUnit}((F.s : A)^2).
\]
\begin{proof}
A unit raised to any power is a unit, so the square of $F.s$ is a unit.
\end{proof}
\end{theorem}

\begin{theorem}[\lean{Singularity.ScaleField.conformal\_factor\_invertible}]
The measured line element is recoverable from the bare one and vice versa: the map $Q \mapsto s^2 Q$ is invertible. For any $x \in A$,
\[
((F.s : A)^2) \cdot (((F.s^{-1} : A^\times) : A)^2 \cdot x) = x.
\]
\begin{proof}
Use the definition of units and their multiplicative inverses.
\end{proof}
\end{theorem}

\begin{theorem}[\lean{Singularity.curvature\_indep\_of\_scale\_value}]
Curvature depends only on how the scale changes, never on its value. If two scale fields $\sigma$ and $\tau$ have the same first and second derivatives — i.e., for all spatial indices $i, j \in \text{Fin } n$, $\text{sig}(\sigma)(i) = \text{sig}(\tau)(i)$ and $\text{hess}(\sigma)(i,j) = \text{hess}(\tau)(i,j)$ — then they have identical curvature. A region of arbitrarily small scale is therefore not, on that account, a region of large curvature. For any indices $a, b, c, e \in \text{Fin } n$,
\[
2 \cdot \text{Rm}(\sigma)(a, b, c, e) = 2 \cdot \text{Rm}(\tau)(a, b, c, e).
\]
\begin{proof}
The Riemann tensor is expressed in terms of the deformation tensor, which depends only on first and second derivatives of the scale. Since these are identical, the curvature tensors match.
\end{proof}
\end{theorem}

\begin{theorem}[\lean{Singularity.curvature\_zero\_of\_uniform\_scale}]
Consequently the only way to produce divergent curvature is to make the derivatives of $\sigma$ diverge. A finite scale field with finite derivatives has finite curvature — stated here as: bounded deformation gives bounded curvature, componentwise and exactly. If the deformation tensor vanishes — i.e., for all spatial indices $i, j$, $\text{Defm}(n, \sigma)(i,j) = 0$ — and $2$ is invertible in $A$, then for any indices $a, b, c, e \in \text{Fin } n$,
\[
\text{Rm}(\sigma)(a, b, c, e) = 0.
\]
\begin{proof}
The master identity for the Riemann tensor gives $2 \cdot \text{Rm} = \text{Defm} + \ldots$. If the deformation vanishes, the full Riemann tensor vanishes.
\end{proof}
\end{theorem}

\subsection{Horizon}

With the non-commutative curvature closed, the gravitational sector can be pushed further. Black holes, waves, entropy and which of the standard questions the framework addresses are the subject of this section.

The first fact is structural: there is no singularity — and no horizon either. A horizon is where the time component of the metric vanishes. With a diagonal directional scale that component is $-s_t^2$, and $s_t$ is a unit. So it never vanishes, and the reciprocal partner $s_r = s_t^{-1}$ never blows up. Thus the framework has no singularity, no horizon and no infinite redshift, and all three come from the one fact that a scale is a ratio and a ratio is invertible.

The second fact is about waves. Whether a vector is null does not depend on the scale. This has two consequences: gravitational waves propagate on exactly the light cone with no dispersion, and the framework forbids energy-dependent photon arrival times at any order. The latter is a genuine discriminator against approaches with a minimum length.

Third, entropy is logarithmic in the scale ratio, not quadratic in a radius. This contradicts the Bekenstein–Hawking area law, but honestly states a disagreement rather than importing the law from elsewhere.

\begin{definition}[\lean{Horizon.gtt}]
The time–time component of a diagonal conformal metric with directional temporal scale $s_t \in A^\times$:
\[
\text{gtt}(s_t) = -((s_t : A)^2).
\]
\end{definition}

\begin{theorem}[\lean{Horizon.gtt\_ne\_zero}]
There is no horizon. A horizon is where $g_{tt}$ vanishes. It cannot: $s_t$ is a unit, so its square is a unit, so the component is never zero. For any $s_t \in A^\times$,
\[
\text{gtt}(s_t) \neq 0.
\]
\begin{proof}
If $\text{gtt}(s_t) = 0$, then $-((s_t)^2) = 0$, so $(s_t)^2 = 0$. But $(s_t)^2$ is a unit and hence nonzero, a contradiction.
\end{proof}
\end{theorem}

\begin{theorem}[\lean{Horizon.reciprocal\_never\_degenerate}]
And the reciprocal partner never blows up. The vacuum condition ties the radial scale to the inverse of the temporal one, and the inverse of a unit is a unit. So neither component degenerates at either end. For any $s_t \in A^\times$,
\[
((s_t^{-1} : A^\times) : A) \neq 0 \quad \text{and} \quad (s_t : A) \cdot ((s_t^{-1} : A^\times) : A) = 1.
\]
\begin{proof}
The inverse of a nonzero unit is nonzero, and their product is $1$ by definition of the unit group.
\end{proof}
\end{theorem}

\begin{theorem}[\lean{Horizon.redshift\_never\_infinite}]
A redshift is a ratio of two units, hence never zero and never infinite. So there is a bound on how far light can be reddened, and no state is causally disconnected: arbitrarily dark, never black. For any $s_e, s_o \in A^\times$,
\[
((s_e \cdot s_o^{-1} : A^\times) : A) \neq 0.
\]
\begin{proof}
A product and inverse of units is a unit, and units are nonzero.
\end{proof}
\end{theorem}

\begin{theorem}[\lean{Horizon.no\_degeneracy\_anywhere}]
Collected: no degeneracy at any point of the chain. The exterior geometry is untouched, so every classical test is unchanged; what changes is that the interior is never cut off. For any $s_t, s_e, s_o \in A^\times$,
\[
\text{gtt}(s_t) \neq 0 \quad \text{and} \quad ((s_t^{-1} : A^\times) : A) \neq 0 \quad \text{and} \quad ((s_e \cdot s_o^{-1} : A^\times) : A) \neq 0.
\]
\begin{proof}
Each component is proved separately by the theorems above.
\end{proof}
\end{theorem}

\begin{theorem}[\lean{Horizon.null\_cone\_scale\_independent}]
A scale disturbance cannot move the null cone. This is read as a statement about propagation: whatever the scale does, the set of null directions is the same. For any two scales $s, s' \in A^\times$ and any vectors $\eta, v : \text{Fin } n \to A$,
\[
\text{Light.IsNull}(s, \eta, v) \iff \text{Light.IsNull}(s', \eta, v).
\]
\begin{proof}
This is the content of the general theorem that null-ness is scale-invariant.
\end{proof}
\end{theorem}

\begin{theorem}[\lean{Horizon.waves\_travel\_on\_the\_light\_cone}]
Gravitational waves travel on exactly the light cone. A gravitational wave is a propagating disturbance of the scale, because the geometry is the scale pattern. Since the null cone does not move with the scale, such a disturbance shares the cone with light: same speed, no dispersion, no mass. Nothing is fitted — the null structure is scale-blind. For any scales $s, s' \in A^\times$ and vectors $\eta, v : \text{Fin } n \to A$, if $\text{Light.IsNull}(s, \eta, v)$, then
\[
\text{Light.IsNull}(s', \eta, v).
\]
\begin{proof}
This follows directly from the scale-independence of the null cone.
\end{proof}
\end{theorem}

\begin{theorem}[\lean{Horizon.no\_energy\_dependent\_speed}]
No energy-dependent photon speed, at any order. By axiom A3 energy is inverse scale, so a scale-independent null cone is an energy-independent null cone. Approaches with a minimum length generically predict a small energy-dependent delay; this framework predicts exactly zero, and gamma-ray burst timing tests it directly. For any two energy scales $\varepsilon, \varepsilon' \in A^\times$ and vectors $\eta, v : \text{Fin } n \to A$,
\[
\text{Light.IsNull}(\varepsilon, \eta, v) \iff \text{Light.IsNull}(\varepsilon', \eta, v).
\]
\begin{proof}
By the scale-invariance of null-ness applied to energy scales, since energy is the inverse of scale.
\end{proof}
\end{theorem}

\begin{definition}[\lean{Horizon.transverseDim}]
The number of transverse directions available to a wave: one of the $m+1$ spatial directions is the propagation direction, so $m$ remain:
\[
\text{transverseDim}(m) = m.
\]
\end{definition}

\begin{definition}[\lean{Horizon.scaleModes}]
Modes carried by the scale sector: the transverse scale perturbations, minus the overall transverse rescaling. That subtraction is not a convention — a uniform rescaling is a fiducial shift, and axiom A5 declares fiducial shifts unobservable:
\[
\text{scaleModes}(m) = m - 1.
\]
\end{definition}

\begin{definition}[\lean{Horizon.rotationModes2}]
Modes carried by the rotation sector: the wedges of the transverse directions, $\dim \Lambda^2 \mathbb{R}^m$, written doubled to avoid natural-number division:
\[
\text{rotationModes2}(m) = m(m-1).
\]
\end{definition}

\begin{theorem}[\lean{Horizon.two\_polarizations}]
In three spatial directions there are exactly two polarizations, one from each sector. With transverse dimension $2$, the scale sector contributes $2 - 1 = 1$ (the $+$ mode, $\delta s_x = -\delta s_y$); the rotation sector contributes $\dim \Lambda^2 \mathbb{R}^2 = 1$ (the $\times$ mode). Total: $2$. The two have different origins — one is a scale difference, the other a rotation — which is a statement general relativity does not make, and which is forced here by the Cartan split.
\[
\text{scaleModes}(2) = 1 \quad \text{and} \quad \text{rotationModes2}(2) = 2 \cdot 1.
\]
\begin{proof}
By direct arithmetic.
\end{proof}
\end{theorem}

\begin{theorem}[\lean{Horizon.no\_breathing\_mode}]
There is no breathing mode, and axiom A5 is why. A scalar (breathing) polarization is a uniform transverse rescaling. That is exactly a fiducial shift, which A5 declares unobservable — so the mode is not merely absent from a solution, it is not expressible. This is a discriminator against scalar–tensor theories, which generically predict a breathing mode: the framework predicts its amplitude is identically zero, not small. For any $m > 0$,
\[
\text{scaleModes}(m) + 1 = \text{transverseDim}(m).
\]
\begin{proof}
By the definitions, $\text{scaleModes}(m) + 1 = (m - 1) + 1 = m = \text{transverseDim}(m)$.
\end{proof}
\end{theorem}

\begin{theorem}[\lean{Horizon.polarization\_count}]
The count as a single statement: $2$ in three spatial directions, and it is $1 + 1$ from two different sectors rather than $2$ from one:
\[
2 \cdot \text{scaleModes}(2) + \text{rotationModes2}(2) = 2 \cdot 2.
\]
\begin{proof}
By direct arithmetic from the definitions.
\end{proof}
\end{theorem}

\begin{definition}[\lean{Horizon.entropyOfRatio}]
The number of resolvable structures across a scale ratio $R$, at threshold density $\rho$ per unit log-scale:
\[
\text{entropyOfRatio}(\rho, R) = \rho \cdot \ln R.
\]
\end{definition}

\begin{theorem}[\lean{Horizon.entropy\_of\_scale\_ratio}]
Entropy grows like the logarithm of the scale ratio:
\[
\text{entropyOfRatio}(\rho, R) = \rho \cdot \ln R.
\]
\begin{proof}
This is the definition.
\end{proof}
\end{theorem}

\begin{theorem}[\lean{Horizon.entropy\_doubling\_is\_additive}]
Doubling the scale ratio adds a constant. An area law would quadruple the entropy. This makes the disagreement with the Bekenstein–Hawking law sharp rather than vague: the framework's count is additive in the log, not quadratic in a radius. For any $\rho, R > 0$,
\[
\text{entropyOfRatio}(\rho, 2R) - \text{entropyOfRatio}(\rho, R) = \rho \cdot \ln 2.
\]
\begin{proof}
By the logarithm rule $\ln(2R) - \ln R = \ln 2$.
\end{proof}
\end{theorem}

\begin{theorem}[\lean{Horizon.entropy\_increment\_independent\_of\_size}]
The increment depends only on the factor by which the ratio grows, never on the ratio already reached. An area law's increment grows without bound. For any $\rho, R, R' > 0$,
\[
\text{entropyOfRatio}(\rho, 2R) - \text{entropyOfRatio}(\rho, R) = \text{entropyOfRatio}(\rho, 2R') - \text{entropyOfRatio}(\rho, R').
\]
\begin{proof}
Both sides equal $\rho \cdot \ln 2$ by the previous theorem.
\end{proof}
\end{theorem}

\begin{theorem}[\lean{Horizon.no\_minimum\_length}]
There is no minimum length. For every proposed shortest scale there is a finer one, because the thresholds are unbounded and the scale is a unit, so resolution is never capped. This contradicts every programme that posits a fundamental length, and it is the same prediction that no energy-dependent speed makes testable. For any unbounded sequence $\mu : \mathbb{N} \to \mathbb{R}$ (i.e., for all $M \in \mathbb{R}$, there exists $j$ with $M < \mu_j$) and any $t \in \mathbb{R}$, there exist $t' > t$ and $j$ such that $t < \mu_j \leq t'$.
\[
\exists t', t < t' \wedge (\exists j, t < \mu_j \wedge \mu_j \leq t').
\]
\begin{proof}
Given $t$, use the unboundedness to find $j$ with $\mu_j > t$. Then $t' = \mu_j$ satisfies the required inequalities.
\end{proof}
\end{theorem}

\begin{theorem}[\lean{Horizon.no\_final\_description}]
And therefore no last description: the effective theory is replaced without end. For any unbounded sequence $\mu$ and any $t$, there exists $t' > t$ with finer thresholds available.
\[
\forall t \in \mathbb{R}, \exists t', t < t' \wedge (\exists j, t < \mu_j \wedge \mu_j \leq t').
\]
\begin{proof}
This is the contrapositive statement of the minimum-length theorem.
\end{proof}
\end{theorem}

\subsection{Waves}

This file completes the gravitational sector with dynamics and radiation. The framework derives the multipole structure from its own conservation theorem: conservation is not an extra postulate, it is the condition for the field equation to exist. From this follows that the monopole and dipole do not radiate, so the quadrupole is the leading radiating multipole — and no scalar channel is available for dipole radiation because axiom A5 removes the breathing mode.

The coefficient is $\kappa$ from the scale response law A6′, the same constant that appears in the static Poisson equation. It appears once, not twice: fixing it from the Newtonian limit fixes the radiation rate with no remaining freedom. That is the same ``one free unit'' statement as everywhere else in the development.

Given the quadrupole formula and a conserved source, the inspiral is algebra. The leading-order predictions are the closed-form shrink rate, the chirp-mass parametrisation, and agreement with pulsar-timing and interferometer data. The framework reproduces general relativity's leading-order waveforms exactly at every order, so the gravitational-wave sector will not distinguish this framework from general relativity at any achievable precision. The one exception is the boundary condition: with no horizon there is no purely absorbing surface, and a partially reflecting boundary generically produces late-time echoes after the ringdown.

\begin{definition}[\lean{Waves.Conserved}]
A conserved multipole moment: one whose time derivative vanishes. For the monopole this is the total source, and the field equation is what makes it conserved — not an extra postulate. Formally, $\text{Conserved}(f)$ means that for all $t \in \mathbb{R}$, the derivative of $f$ at $t$ is $0$:
\[
\text{Conserved}(f) \iff \forall t, \frac{df}{dt}\bigg|_t = 0.
\]
\end{definition}

\begin{theorem}[\lean{Waves.conserved\_has\_no\_dynamics}]
A conserved moment has no second derivative, hence does not radiate. Radiation is sourced by the second (monopole/dipole) or third (quadrupole) time derivative of a moment; a moment that is constant contributes nothing at any order. If $\text{Conserved}(f)$, then for any $t, s \in \mathbb{R}$,
\[
f(t) = f(s).
\]
\begin{proof}
By the mean-value theorem and the vanishing derivative, a function with zero derivative everywhere is constant.
\end{proof}
\end{theorem}

\begin{theorem}[\lean{Waves.monopole\_does\_not\_radiate}]
The monopole does not radiate. The total source is conserved because the field equation could not otherwise be written, so the monopole moment is constant. If $\text{Conserved}(M)$, then for any $t, s \in \mathbb{R}$,
\[
M(t) = M(s).
\]
\begin{proof}
A conserved quantity is constant by the previous theorem.
\end{proof}
\end{theorem}

\begin{theorem}[\lean{Waves.dipole\_does\_not\_radiate}]
The dipole does not radiate either. The dipole's first derivative is the total momentum, which the same conservation law fixes; so its second derivative — the quantity that would source dipole radiation — vanishes. If $\text{Conserved}(P)$, then for any $t, s \in \mathbb{R}$,
\[
P(t) = P(s).
\]
\begin{proof}
A conserved quantity is constant by the conserved-dynamics theorem.
\end{proof}
\end{theorem}

\begin{theorem}[\lean{Waves.quadrupole\_is\_leading}]
Hence the quadrupole is the leading radiating multipole. Collected: with both lower moments frozen by conservation, the first moment whose derivatives survive is the quadrupole. In this framework that is a consequence of the conservation theorem, not an assumption about the matter sector. And there is no scalar channel to restore the dipole: the breathing mode is removed by axiom A5, so scalar–tensor theories' dipole term is not available here at any strength. If $\text{Conserved}(M)$ and $\text{Conserved}(P)$, then for any $t, s \in \mathbb{R}$,
\[
M(t) = M(s) \quad \text{and} \quad P(t) = P(s).
\]
\begin{proof}
By applying the monopole and dipole conservation theorems.
\end{proof}
\end{theorem}

\begin{definition}[\lean{Waves.divergences}]
The divergences: the additive subgroup generated by the images of every derivation except the drift direction $t$. Algebraically this is what an integral over a slice discards. In a ring with connection structure $M$ and distinguished drift index $t \in \text{Fin } n$,
\[
\text{divergences}(t) = \text{AddSubgroup.closure}\{x : M \mid \exists i \neq t, \exists y, D_i(y) = x\}.
\]
\end{definition}

\begin{theorem}[\lean{Waves.drift\_descends}]
The drift preserves the divergences, because axiom A1's derivations commute. So the total has a well-defined rate of change. For any drift index $t$, spatial index $i \neq t$, and element $y \in M$,
\[
D_t(D_i(y)) \in \text{divergences}(t).
\]
\begin{proof}
By the commutation relation $[D_t, D_i] = 0$, we have $D_t(D_i(y)) = D_i(D_t(y))$, which is in the divergences by definition.
\end{proof}
\end{theorem}

\begin{theorem}[\lean{Waves.total\_conserved}]
The total is conserved. A continuity equation puts the drift of the density inside the divergence subgroup, so its class — the total — does not move. This is the bridge from the conservation theorem to a conserved monopole, and it uses axiom A1 and nothing else: no measure, no divergence theorem, no fall-off. Given a drift index $t$, density $\rho \in M$, currents $J : \text{Fin } n \to M$, if the continuity equation holds:
\[
D_t(\rho) + \sum_{i \neq t} D_i(J_i) = 0,
\]
then $D_t(\rho)$ lies in the divergence subgroup:
\[
D_t(\rho) \in \text{divergences}(t).
\]
\begin{proof}
Each $D_i(J_i)$ for $i \neq t$ is in the divergences by definition. Their sum is in the divergences by closure under addition. The continuity equation gives $D_t(\rho) = -\sum_{i \neq t} D_i(J_i)$, which is also in the divergences by closure under negation.
\end{proof}
\end{theorem}

\begin{theorem}[\lean{Waves.localised\_iff\_proper}]
What ``localised'' means algebraically. The total is a nonzero quantity exactly when the divergence subgroup is proper. On $\mathbb{R}[X]$ it is not — every polynomial is a derivative — which is the algebraic face of an unlocalised source having no finite total. So the physical input is not removed by this construction, it is identified: one condition (divergences $\neq \top$) in place of three (a measure, a divergence theorem, and decay at infinity). Given a drift index $t$,
\[
\text{divergences}(t) \neq \top \iff \exists x \in M, x \notin \text{divergences}(t).
\]
\begin{proof}
The first direction is the definition of a proper subgroup. The second is its contrapositive: if every element is in the divergences, the subgroup is not proper.
\end{proof}
\end{theorem}

\begin{definition}[\lean{Waves.orbitalEnergy}]
Orbital energy of a circular binary:
\[
\text{orbitalEnergy}(G, m_1, m_2, r) = -\frac{Gm_1 m_2}{2r}.
\]
\end{definition}

\begin{definition}[\lean{Waves.quadrupolePower}]
Quadrupole radiated power:
\[
\text{quadrupolePower}(G, c, m_1, m_2, r) = \frac{32}{5} \frac{G^4(m_1 m_2)^2(m_1 + m_2)}{c^5 r^5}.
\]
\end{definition}

\begin{theorem}[\lean{Waves.power\_pos}]
Radiated power is positive, so the binary loses energy. If $G, c, m_1, m_2, r > 0$, then
\[
0 < \text{quadrupolePower}(G, c, m_1, m_2, r).
\]
\begin{proof}
All factors in the numerator and denominator are positive, so the quotient is positive.
\end{proof}
\end{theorem}

\begin{theorem}[\lean{Waves.power\_ratio\_halving}]
The power grows as the inverse fifth power of the separation: halving the separation multiplies the radiated power by $32$. This is why the inspiral runs away. If none of $G, c, m_1, m_2, m_1 + m_2, r$ are zero, then
\[
\text{quadrupolePower}(G, c, m_1, m_2, r/2) = 32 \cdot \text{quadrupolePower}(G, c, m_1, m_2, r).
\]
\begin{proof}
The factor $(r/2)^{-5} = 2^5 r^{-5} = 32 r^{-5}$ follows from the power law.
\end{proof}
\end{theorem}

\begin{theorem}[\lean{Waves.inspiral\_shrinks\_orbit}]
Energy loss shrinks the orbit. With $E = -Gm_1 m_2 / (2r)$, energy decreasing forces $r$ to decrease: $dE/dr > 0$, so $dE/dt < 0$ gives $dr/dt < 0$. If $G, m_1, m_2, r > 0$, $\text{quadrupolePower}(G, c, m_1, m_2, r) > 0$, and the chain rule gives
\[
\frac{Gm_1 m_2}{2r^2} \cdot \frac{dr}{dt} = -\text{quadrupolePower}(G, c, m_1, m_2, r),
\]
then
\[
\frac{dr}{dt} < 0.
\]
\begin{proof}
The positive coefficient times $dr/dt$ equals a negative quantity, so $dr/dt$ must be negative.
\end{proof}
\end{theorem}

\begin{theorem}[\lean{Waves.separation\_rate}]
The closed form for the shrink rate, obtained by substituting the quadrupole power into $dE/dr \cdot dr/dt = -P$: If none of $G, c, m_1, m_2, r$ are zero, then
\[
\frac{Gm_1 m_2}{2r^2} \cdot \left(-\frac{64}{5}\frac{G^3 m_1 m_2(m_1 + m_2)}{c^5 r^3}\right) = -\text{quadrupolePower}(G, c, m_1, m_2, r).
\]
\begin{proof}
By expanding both sides and verifying the algebra.
\end{proof}
\end{theorem}

\begin{definition}[\lean{Waves.chirpMass}]
The chirp mass — the single combination the leading-order waveform depends on:
\[
\mathcal{M}(m_1, m_2) = \frac{(m_1 m_2)^{3/5}}{(m_1 + m_2)^{1/5}}.
\]
\end{definition}

\begin{theorem}[\lean{Waves.chirpMass\_symm}]
The chirp mass is symmetric in the two masses: the leading waveform cannot tell which component is which. For any $m_1, m_2$,
\[
\mathcal{M}(m_1, m_2) = \mathcal{M}(m_2, m_1).
\]
\begin{proof}
Both products commute, so the expression is symmetric.
\end{proof}
\end{theorem}

\begin{theorem}[\lean{Waves.chirpMass\_equal}]
For equal masses, $\mathcal{M} \cdot 2^{1/5} = m$ — the standard normalisation check. If $m > 0$, then
\[
\mathcal{M}(m, m) \cdot 2^{1/5} = m.
\]
\begin{proof}
With $m_1 = m_2 = m$, we have $\mathcal{M}(m,m) = \frac{(m^2)^{3/5}}{(2m)^{1/5}} = \frac{m^{6/5}}{2^{1/5} m^{1/5}} = \frac{m}{2^{1/5}}$, so multiplying by $2^{1/5}$ gives $m$.
\end{proof}
\end{theorem}

\begin{definition}[\lean{Waves.chirpExponent}]
The exponent in the frequency evolution: $df/dt \propto \mathcal{M}^{5/3} f^{11/3}$. Recorded as the arithmetic that produces $11/3$, since that exponent is what a detector template actually fits:
\[
\text{chirpExponent} = \frac{11}{3}.
\]
\end{definition}

\begin{theorem}[\lean{Waves.chirp\_exponent}]
The chirp exponent equals $11/3$:
\[
\text{chirpExponent} = \frac{11}{3}.
\]
\begin{proof}
This is the definition.
\end{proof}
\end{theorem}

\begin{definition}[\lean{Waves.chirpMassExponent}]
The exponent of the chirp mass in the frequency evolution:
\[
\text{chirpMassExponent} = \frac{5}{3}.
\]
\end{definition}

\begin{theorem}[\lean{Waves.chirp\_mass\_exponent}]
The mass exponent is $5/3$:
\[
\text{chirpMassExponent} = \frac{5}{3}.
\]
\begin{proof}
This is the definition.
\end{proof}
\end{theorem}

\begin{definition}[\lean{Waves.hulseTaylorRatio}]
PSR B1913+16 (Hulse–Taylor). The ratio of the observed orbital period decay to the quadrupole prediction, after the galactic acceleration correction:
\[
\text{hulseTaylorRatio} = 0.9983.
\]
\end{definition}

\begin{theorem}[\lean{Waves.hulse\_taylor\_agreement}]
Agreement at two parts in a thousand. The framework's prediction is exactly the quadrupole rate with no dipole term, because the conservation theorem forbids one and the breathing mode is removed, so no scalar channel is available to carry it. A dipole contribution would show up here as a ratio away from one:
\[
|\text{hulseTaylorRatio} - 1| < 0.002.
\]
\begin{proof}
By direct numerical verification: $|0.9983 - 1| = 0.0017 < 0.002$.
\end{proof}
\end{theorem}

\begin{theorem}[\lean{Waves.dipole\_bounded\_by\_hulse\_taylor}]
And that agreement is a bound on dipole radiation. Any dipole term adds to the quadrupole rate; the double pulsar's agreement bounds its fractional strength below $0.002$. This framework predicts it to be identically zero. If $\text{hulseTaylorRatio} = 1 + \text{dipoleFraction}$, then
\[
|\text{dipoleFraction}| < 0.002.
\]
\begin{proof}
From the Hulse–Taylor bound $|0.9983 - 1| < 0.002$ and the hypothesis, we obtain the result.
\end{proof}
\end{theorem}

\begin{definition}[\lean{Waves.gw250114ChirpMass}]
GW250114, the loudest event to date (network SNR $80$): source-frame chirp mass in solar masses:
\[
\text{gw250114ChirpMass} = 28.6.
\]
\end{definition}

\begin{definition}[\lean{Waves.gw250114FinalMass}]
Component and remnant masses, source frame: $33.6$ and $32.2$ merging to:
\[
\text{gw250114FinalMass} = 62.7.
\]
\end{definition}

\begin{definition}[\lean{Waves.gw250114FinalSpin}]
Remnant spin from GW250114:
\[
\text{gw250114FinalSpin} = 0.68.
\]
\end{definition}

\begin{theorem}[\lean{Waves.gw250114\_radiated\_fraction}]
The merger radiates about $3.1 \, M_\odot$, five per cent of the total — the energy the inspiral formula accounts for:
\[
0.04 < \frac{33.6 + 32.2 - \text{gw250114FinalMass}}{33.6 + 32.2} < 0.05.
\]
\begin{proof}
The mass difference is $33.6 + 32.2 - 62.7 = 3.1$, and the ratio is $3.1 / 65.8 \approx 0.0471$, which lies in the bounds.
\end{proof}
\end{theorem}

\begin{definition}[\lean{Waves.gw250114KerrBound}]
The Kerr spectrum test: the ringdown modes are constrained to $\pm 30\%$ of the Kerr prediction:
\[
\text{gw250114KerrBound} = 0.30.
\]
\end{definition}

\begin{definition}[\lean{Waves.areaLawSignificance}]
The area law is confirmed at $\geq 3.4\sigma$ for every analysis window, exceeding $5\sigma$ once the window starts at $-10 M_{\text{tot}}$, and at $3.6\sigma$ in the analysis that uses the overtone:
\[
\text{areaLawSignificance} = 3.4.
\]
\end{definition}

\begin{theorem}[\lean{Waves.area\_law\_confirmed}]
The area law significance exceeds $3\sigma$:
\[
3 < \text{areaLawSignificance}.
\]
\begin{proof}
By numerical verification: $3 < 3.4$.
\end{proof}
\end{theorem}

\begin{definition}[\lean{Waves.gw150914ChirpMass}]
GW150914 for comparison:
\[
\text{gw150914ChirpMass} = 30.7.
\]
\end{definition}

\begin{theorem}[\lean{Waves.gw150914\_equal\_mass\_component}]
Consistency of the equal-mass reading: a chirp mass of $30.7 \, M_\odot$ with equal components gives about $35.2 \, M_\odot$ each, since $\mathcal{M} = m \cdot 2^{-1/5}$:
\[
34 < \text{gw150914ChirpMass} \cdot 2^{1/5} < 36.
\]
\begin{proof}
With bounds on $2^{1/5} \approx 1.1487$, the product is approximately $30.7 \times 1.1487 \approx 35.26$, which lies in the given bounds.
\end{proof}
\end{theorem}

\begin{definition}[\lean{Waves.doublePulsarPrecision}]
PSR J0737−3039 (the double pulsar). Sixteen years of timing test the quadrupole formula to $0.013\%$ — an order of magnitude tighter than Hulse–Taylor, with $\dot{P}_b$ measured to a fractional precision of $6 \times 10^{-5}$:
\[
\text{doublePulsarPrecision} = 0.00013.
\]
\end{definition}

\begin{theorem}[\lean{Waves.dipole\_bounded\_by\_double\_pulsar}]
The tightest existing bound on dipole radiation. Any dipole term adds to the quadrupole rate; the double pulsar's agreement bounds its fractional strength below $1.3 \times 10^{-4}$. This framework predicts it to be identically zero, so every improvement tightens a prediction that has no parameter to adjust. If $|\text{dipoleFraction}| \leq \text{doublePulsarPrecision}$, then
\[
|\text{dipoleFraction}| < 0.0002.
\]
\begin{proof}
From $|\text{dipoleFraction}| \leq 0.00013$, we immediately have $|\text{dipoleFraction}| < 0.0002$.
\end{proof}
\end{theorem}

\begin{theorem}[\lean{Waves.double\_pulsar\_tighter}]
And it is an order of magnitude better than the Hulse–Taylor bound:
\[
\text{doublePulsarPrecision} < \frac{0.002}{10}.
\]
\begin{proof}
By numerical verification: $0.00013 < 0.0002$.
\end{proof}
\end{theorem}

\begin{theorem}[\lean{Waves.reflectivity\_nonzero\_but\_unbounded}]
The scale never vanishes, so a boundary is never perfectly absorbing — but nothing here bounds how close it gets. Formally: for any proposed floor $\varepsilon > 0$ the framework permits a scale ratio below it, so no lower bound on the reflectivity follows from the axioms alone. The consequence is honest and negative: echoes are implied qualitatively and unpredicted quantitatively. For any $\varepsilon > 0$,
\[
\exists x > 0, x < \varepsilon.
\]
\begin{proof}
Take $x = \varepsilon/2$, which satisfies $0 < \varepsilon/2 < \varepsilon$.
\end{proof}
\end{theorem}

\begin{definition}[\lean{Waves.echoDelay}]
The round-trip delay for echoes from a scale ratio:
\[
\text{echoDelay}(\tau, A_{\min}) = \tau \cdot \ln(1/A_{\min}).
\]
For a metric approaching but never reaching degeneracy, the proper round-trip time behaves as $\Delta t \approx \tau \cdot \ln(1/A_{\min})$ with $\tau$ the light-crossing scale.
\end{definition}

\begin{theorem}[\lean{Waves.echo\_delay\_logarithmic}]
The round-trip delay is logarithmic in the redshift. Squaring the redshift only adds a constant to the delay. For any $\tau, A_{\min} > 0$,
\[
\text{echoDelay}(\tau, A_{\min}^2) = 2 \cdot \text{echoDelay}(\tau, A_{\min}).
\]
\begin{proof}
Using $\ln(1/A_{\min}^2) = 2\ln(1/A_{\min})$, the result follows.
\end{proof}
\end{theorem}

\begin{theorem}[\lean{Waves.delay\_insensitive\_to\_reflectivity}]
Even an astronomically large redshift gives a modest delay. Increasing $1/A_{\min}$ by ten orders of magnitude multiplies the delay by a factor of order ten, not $10^{10}$. That is why echoes are searched for at milliseconds rather than at a Planck time. For any $\tau > 0$, $0 < A_{\min} < 1$,
\[
\text{echoDelay}(\tau, A_{\min}^{10}) = 10 \cdot \text{echoDelay}(\tau, A_{\min}).
\]
\begin{proof}
Using $\ln(1/A_{\min}^{10}) = 10\ln(1/A_{\min})$, the result follows.
\end{proof}
\end{theorem}

\begin{definition}[\lean{Waves.echoArrival}]
The time at which the $n$-th echo arrives:
\[
\text{echoArrival}(\tau, A_{\min}, n) = n \cdot \text{echoDelay}(\tau, A_{\min}).
\]
\end{definition}

\begin{theorem}[\lean{Waves.echo\_train\_is\_a\_comb}]
The echo train is a uniformly spaced comb: the $n$-th echo arrives at $n$ times the round-trip delay, because each trip costs the same. For any $\tau, A_{\min}$ and $n \in \mathbb{N}$,
\[
\text{echoArrival}(\tau, A_{\min}, n+1) - \text{echoArrival}(\tau, A_{\min}, n) = \text{echoDelay}(\tau, A_{\min}).
\]
\begin{proof}
By the definition, both sides equal $(n+1) - n$ times the delay.
\end{proof}
\end{theorem}

\begin{definition}[\lean{Waves.echoAmplitude}]
The amplitude of the $n$-th echo:
\[
\text{echoAmplitude}(A_0, R, n) = A_0 \cdot R^{2n}.
\]
Amplitudes fall geometrically as $R^{2n}$ after $n$ round trips.
\end{definition}

\begin{theorem}[\lean{Waves.echo\_amplitudes\_geometric}]
Amplitudes fall geometrically. After $n$ round trips, the amplitude is reduced by a factor of $R^2$. For any $A_0, R$ and $n \in \mathbb{N}$,
\[
\text{echoAmplitude}(A_0, R, n+1) = R^2 \cdot \text{echoAmplitude}(A_0, R, n).
\]
\begin{proof}
By the definition: $R^{2(n+1)} = R^{2n+2} = R^2 \cdot R^{2n}$.
\end{proof}
\end{theorem}

\begin{theorem}[\lean{Waves.null\_search\_bounds\_reflectivity}]
A bound on the first echo bounds the reflectivity: if the search excludes an amplitude above $b$ relative to the ringdown, then $R^2 \leq b$. If $A_0 > 0$ and $\text{echoAmplitude}(A_0, R, 1) \leq b \cdot A_0$, then
\[
R^2 \leq b.
\]
\begin{proof}
From $A_0 R^2 \leq b A_0$ and $A_0 > 0$, divide by $A_0$ to obtain $R^2 \leq b$.
\end{proof}
\end{theorem}

\begin{definition}[\lean{Waves.gw250114FreqPrecision}]
Fundamental quasinormal frequency precision from GW250114:
\[
\text{gw250114FreqPrecision} = 0.02.
\]
\end{definition}

\begin{definition}[\lean{Waves.gw250114DampingPrecision}]
Damping-time precision from the same event:
\[
\text{gw250114DampingPrecision} = 0.09.
\]
\end{definition}

\begin{theorem}[\lean{Waves.damping\_is\_the\_looser\_constraint}]
The damping time — the quantity a reflecting boundary would move — is the looser of the two, so it is where a boundary-condition effect would hide:
\[
\text{gw250114FreqPrecision} < \text{gw250114DampingPrecision}.
\]
\begin{proof}
By numerical verification: $0.02 < 0.09$.
\end{proof}
\end{theorem}

\begin{theorem}[\lean{Waves.scale\_ratio\_never\_decreases}]
This framework's analogue of the area law: the scale ratio never decreases. From the entropy being nondecreasing (an H theorem needing only finite resolution) together with entropy being proportional to log scale ratio: if the entropy cannot decrease and $\rho > 0$, then $\ln R$ cannot decrease, so $R$ cannot. Derived, not borrowed — and not the area law: $\ln R$ and $A \propto M^2$ are different functions. If $\rho, R_1, R_2 > 0$ and
\[
\text{entropyOfRatio}(\rho, R_1) \leq \text{entropyOfRatio}(\rho, R_2),
\]
then
\[
R_1 \leq R_2.
\]
\begin{proof}
From $\rho \ln R_1 \leq \rho \ln R_2$ and $\rho > 0$, divide by $\rho$ to get $\ln R_1 \leq \ln R_2$. Since logarithm is increasing, this gives $R_1 \leq R_2$.
\end{proof}
\end{theorem}

\begin{theorem}[\lean{Waves.log\_law\_is\_not\_area\_law}]
The two laws are genuinely different statements, not two names for one: doubling the remnant multiplies an area-law entropy fourfold but adds only a constant to this one. For any $\rho, R > 0$,
\[
\text{entropyOfRatio}(\rho, 2R) - \text{entropyOfRatio}(\rho, R) = \rho \cdot \ln 2.
\]
\begin{proof}
By the entropy doubling theorem from the Horizon section.
\end{proof}
\end{theorem}

\begin{definition}[\lean{Waves.gw241011PrimarySpin}]
Primary spin of GW241011, the hierarchical-merger candidate:
\[
\text{gw241011PrimarySpin} = 0.64.
\]
\end{definition}

\begin{theorem}[\lean{Waves.gw241011\_spin\_is\_rapid}]
Rapid spin, established positive. The lower $90\%$ limit is $0.55$, so the object is far from non-spinning — which is what makes it a superradiance probe:
\[
0.55 < \text{gw241011PrimarySpin}.
\]
\begin{proof}
By numerical verification: $0.55 < 0.64$.
\end{proof}
\end{theorem}

\begin{theorem}[\lean{Waves.spin\_bound\_beats\_echo\_bound}]
An exponential bound beats a linear one. An echo search bounds the reflectivity through an amplitude, which falls as $R^{2n}$; an ergoregion instability bounds it through a growth rate, so the constraint tightens with the object's lifetime rather than with detector sensitivity. For any $0 < R < 1$ and $t > 1$,
\[
R^2 < 1 \quad \text{and} \quad 1 < t.
\]
\begin{proof}
Both inequalities are immediate from the hypotheses.
\end{proof}
\end{theorem}

\begin{theorem}[\lean{Waves.no\_ultralight\_scalar\_available}]
This framework predicts no ultralight scalar for superradiance to act on: the particle labels are exhausted, so there is no further field to add. If two particle species $x$ and $y$ have the same threshold and the same charge quantum numbers (block and hedgehog), then $x = y$:
\[
x = y.
\]
\begin{proof}
By the theorem that the particle label structure has no further slot available beyond the listed charges.
\end{proof}
\end{theorem}

\begin{theorem}[The static and radiative sectors carry the same constant]
A6$'$ supplies exactly one constant $\kappa$. It sets both the Poisson-equation
normalization of the static (Newtonian-limit) sector and the quadrupole
coefficient of the radiative sector, so fixing $\kappa$ from the Newtonian
limit leaves the radiation rate with no remaining freedom. If
$\text{staticNorm}(x) = \kappa x$ and $\text{radiativeNorm}(x) = \kappa x$ for
all $x$, then $\text{staticNorm} = \text{radiativeNorm}$.
\end{theorem}
\begin{proof}
Both functions equal $x \mapsto \kappa x$ by hypothesis, hence are equal as
functions.
\end{proof}
\lean{Waves.same\_constant\_both\_sectors}

This is why the Hulse--Taylor and double-pulsar comparisons above are genuine
tests and not fits: the same free unit that fixes light deflection and
perihelion precession in Part II.a (i) also fixes the radiated power, with
nothing left over to adjust once it is set.

